\noindent People have been emigrating since the beginning of mankind, often starting their life anew in a place they knew little about. It has long been recognised that emigrants' chances to succeed in their new environment depend extensively on their networks \citep{munshi2003networks}. Networks have been found to route workers to jobs and firms \citep{calvo2004effects, beaman2012social}, to carry insurance \citep{townsend1994risk, fafchamps2007formation, ambrus2014consumption}, even determining who migrates in the first place \citep{munshi2016networks, morten2019temporary}. \\

\noindent Previous literature has mostly focused on realized networks, \textit{i.e.} the case where an emigrant has access to a network stock upon arrival. This thesis does not assume networks to be realized on arrival, but seeks to study new networks on the point of formation. I will dub this specific type of network\footnote{Namely, in formation and belonging to a new emigrant.} \textit{"newcomer's network"} throughout. \\

\noindent Ties seem to form disproportionately between people who resemble each other, a social science regularity named \textit{homophily} \citep{mcpherson2001birds}. The literature does not agree on why this occurs. However, having a clear understanding of homophily and its drivers is of key importance. A homophilous network may become a problem, especially for newcomers, who do not yet have the ropes of their new environment. Indeed, if newcomers are tied mostly to other newcomers, the people they can turn to are short of the same things they are \citep{beaman2012social}. \\

\noindent In order to address network homophily, one must disentangle its drivers, as each implies a different solution. For example, if networks are homophilous because individuals only ever meet their own, institutions can rearrange who meets whom. If they meet everyone and still prefer their own, rearranging does little by itself \citep{lowe2021types}. To understand which mechanisms lead to homophily in networks of newcomers is thus the purpose of this thesis. I attempt to do so in a setting where competing explanations can be told apart, namely a garment production zone in southern Ethiopia. The data concerns rural immigrant workers that have just been hired and moved to the city for the occasion.  \\

\noindent In the thesis, I identify four possible ways a newcomer's network may end up homophilous. The newcomer may have a realized network, because she arrived with fellow emigrants or because she is joining family (\textit{importation}). She may have met only her own (\textit{opportunity}), have met everybody and preferred her own (\textit{taste}), or may have tried to befriend people outside her group and been turned down (\textit{exclusion}). I write a small model in which each of the four leaves a different trace in the data, and I attempt to retrieve the effects. This design is possible thanks to a very particular dataset. The survey I use, collected for a forthcoming project by Agarwal, Grosset-Touba and Wu, allows me to build network data which holds three features rarely found together. Firstly, it catches ties at the point of formation, which is useful to isolate importation. Secondly, the set out of which people choose their ties is recorded, allowing me to differentiate who she met from who she chose. Thirdly, nominations are directed, and for a good share of pairs I observe both directions. This is very relevant for exclusion. \\

\noindent I find that the homophily in these networks was mostly imported. This is due to people still relying heavily on their families. For what concerns new relationships, sharing a village of origin is not a relevant predictor. Preference for one's own is absent at the village level and weak at the district level. The more relevant dimension is proximity. Being put on the same production line raises the chance of a relationship by \pnChanOppUnitPp{} percentage points, and two workers with no language in common tie at \pnChanOppKappa{} of the rate of two who can talk to each other. Minority workers, defined by language spoken and region of birth, are named somewhat less by the majority. Finally, sharing an origin does not make a relationship deeper. It does appear to change what the relationship is used for, but mostly because co-villagers tend to be relatives. \\


\noindent This thesis thus makes two contributions. The first is the documentation of network dynamics in this specific and very relevant setting.  In developing economies, foreign firms seeking lower labour costs play an active role in rural-to-urban migration, drawing workers out of agriculture into factory employment \citep{barrett2019made}. The setting of the study, Ethiopia's Hawassa Industrial Park, attracted tens of thousands of rural workers within a few years of opening \citep{mains2021ethiopian}. The same pattern recurs across South-East Asia \citep{heath2018why}. Understanding the beginning of these people's urban life, and how their networks are shaped, can hopefully shed light on the challenges they face and help to improve their long-term outcomes. The second is one of design, \textit{i.e.} how to make the best use of this dataset. \\

\noindent The thesis proceeds as follows. Section~\ref{sec:cf} builds the conceptual framework and reviews the related literature. Section~\ref{sec:setting} describes the setting and data sources used. Section~\ref{sec:model} presents a toy model. Section~\ref{sec:strategy} sets out the empirical strategy, Section~\ref{sec:results} reports the results, and Section~\ref{sec:conclusion} concludes.
